% TC36 — แอดมินจัดการภาพหน้าแรก / หน้านิสิต
\begin{testscenario}{36}{3}
  \tcmeta{TC36}{CS 2568/01}{ระบบเว็บแอปพลิเคชัน ``NUCARE''}{นายสมชาย ใจดี}{วันที่ 7 พฤศจิกายน พ.ศ.\ 2569}{1.0}
  \tcname{การจัดการภาพประชาสัมพันธ์โดยแอดมิน}
  \tcpre{บัญชีแอดมินตาม \tccode{data.md} · รัน \tccode{npm run test:unit} ที่ \tccode{server/} ให้ผ่านเคส SM-*}
  \tcdesc{ตรวจการตรวจสอบชนิด/ขนาดรูป กัน path traversal ของ URL สาธารณะ และการเพิ่มลบผ่าน API ของแอดมิน}

  \tccase{1}{รับเฉพาะ slot และชนิดรูปที่อนุญาต}{%
    \begin{tcsteps}
      \item เรียก \tccode{isSiteMediaSlot} / \tccode{parseImageDataUrl}
      \item ส่ง data URL PNG 1x1 และชนิด GIF / พาธธรรมดา
    \end{tcsteps}}{%
    PNG ผ่าน · GIF และพาธไม่ใช่ data URL ถูกปฏิเสธ}{%
    SM-01/SM-02 ผ่าน}{Pass}

  \tccase{2}{จำกัดขนาดและ sanitize URL}{%
    \begin{tcsteps}
      \item สร้าง buffer ใหญ่กว่า 2~MiB แล้ว parse
      \item เรียก \tccode{toSiteMediaPublicUrl} ด้วยชื่อที่มี \tccode{..}
    \end{tcsteps}}{%
    ไฟล์ใหญ่ถูกปฏิเสธ · URL ใช้เฉพาะ basename}{%
    SM-03/SM-04 ผ่าน}{Pass}

  \tccase{3}{แอดมินเพิ่ม/ลบแล้วหน้าเว็บดึง API}{%
    \begin{tcsteps}
      \item เข้า \tccode{/admin/site-media} เพิ่มภาพช่อง HOME\_HERO / HOME\_GALLERY / STUDENT\_ACTIVITY
      \item เปิดหน้าแรกและหน้านิสิต · ลบรายการหนึ่งแล้วรีเฟรช
    \end{tcsteps}}{%
    แสดงภาพจาก \tccode{/api/uploads/site-media/...} · ลบแล้วไม่แสดง}{%
    ตรวจด้วยมือหลัง deploy โลคัล}{Pass}
\end{testscenario}
